%% titre: Méthodes — Étudier la coercivité d'une fonction
%% matiere: Optimisation
%% niveau: Licence 3
%% annee: 2026-2027
%% ordre: 7
%% description: Vérifier la coercivité ou la non-coercivité d'une fonction par l'étude de sa limite à l'infini.

\section{Démarche}

\begin{methode}[Montrer qu'une fonction $f:\R^n\to\R$ est coercive]
\begin{enumerate}[label=\textbf{Étape \arabic*.}]
    \item \textbf{Rappeler la définition} : une fonction $f$ est coercive si
    \[
    \lim_{\|x\|\to+\infty}f(x)=+\infty.
    \]

    \item \textbf{Développer l'expression} : identifier les termes dominants de $f(x)$ lorsque $\|x\|\to+\infty$.

    \item \textbf{Étudier la limite} : montrer que, quelle que soit la direction suivie pour s'éloigner de l'origine,
    \[
    f(x)\longrightarrow+\infty.
    \]

    \item \textbf{Conclure} : si cette limite vaut $+\infty$, alors
    \[
    \boxed{f\text{ est coercive}.}
    \]
\end{enumerate}
\end{methode}

\begin{attention}
\begin{itemize}
    \item Il faut étudier le comportement de $f$ dans toutes les directions lorsque $\|x\|\to+\infty$.
    \item Il ne suffit pas de vérifier que $f(x)\to+\infty$ le long d'une seule direction.
    \item Pour montrer que $f$ n'est pas coercive, il suffit de trouver une seule suite $(x_k)$ telle que
    \[
    \|x_k\|\to+\infty
    \quad\text{et}\quad
    f(x_k)\not\to+\infty.
    \]
    \item Une fonction coercive n'est pas nécessairement convexe.
\end{itemize}
\end{attention}

\begin{exemple}[Application : montrer la coercivité]
Considérons la fonction
\[
f(x,y)=x^2+2y^2-3x+4y+5.
\]

On complète les carrés :
\[
\begin{aligned}
f(x,y)
&=\left(x-\frac32\right)^2
+2(y+1)^2+\frac34.
\end{aligned}
\]

Lorsque $\|(x,y)\|\to+\infty$, au moins une des coordonnées $x$ ou $y$ devient non bornée. Les termes quadratiques dominent les termes linéaires.

Plus rigoureusement, en utilisant
\[
-3x\geq -\frac12x^2-\frac92,
\qquad
4y\geq -y^2-4,
\]
on obtient
\[
f(x,y)\geq
\frac12x^2+y^2-\frac{17}{2}.
\]

Or,
\[
\frac12x^2+y^2
\geq\frac12(x^2+y^2).
\]

Ainsi,
\[
f(x,y)\geq
\frac12\|(x,y)\|^2-\frac{17}{2}.
\]

Par conséquent,
\[
\lim_{\|(x,y)\|\to+\infty}f(x,y)=+\infty.
\]

Donc
\[
\boxed{f\text{ est coercive}.}
\]
\end{exemple}

\section{Le cas de la non-coercivité}

\begin{methode}[Montrer qu'une fonction n'est pas coercive]
\begin{enumerate}[label=\textbf{Étape \arabic*.}]
    \item \textbf{Choisir une direction} : construire une suite ou une famille de points $x_t$ telle que
    \[
    \|x_t\|\longrightarrow+\infty.
    \]

    \item \textbf{Calculer les valeurs de la fonction} : déterminer le comportement de
    \[
    f(x_t).
    \]

    \item \textbf{Vérifier la limite} : montrer que
    \[
    f(x_t)\not\longrightarrow+\infty.
    \]

    \item \textbf{Conclure} : l'existence d'une telle famille suffit pour affirmer que
    \[
    \boxed{f\text{ n'est pas coercive}.}
    \]
\end{enumerate}
\end{methode}

\begin{exemple}[Application : montrer la non-coercivité]
Considérons
\[
f(x,y)=x^3+y^3-3x-12y+20.
\]

Choisissons la famille
\[
P_t=(-t,0),\qquad t\geq0.
\]

On a
\[
\|P_t\|=t\longrightarrow+\infty.
\]

D'autre part,
\[
\begin{aligned}
f(-t,0)
&=(-t)^3+0-3(-t)-0+20\\
&=-t^3+3t+20.
\end{aligned}
\]

Ainsi,
\[
\lim_{t\to+\infty}f(-t,0)=-\infty.
\]

Par conséquent,
\[
\boxed{f\text{ n'est pas coercive}.}
\]
\end{exemple}

\begin{attention}
La coercivité est une propriété globale qui concerne toutes les directions de fuite vers l'infini. En revanche, la non-coercivité peut être établie à partir d'une seule direction.
\end{attention}

\section{Lien avec l'existence d'un minimum global}

\begin{methode}[Théorème d'existence d'un minimum global]
Soit $f:\R^n\to\R$ une fonction continue et coercive. Alors il existe $x^*\in\R^n$ tel que
\[
f(x^*)=\min_{x\in\R^n}f(x).
\]

La continuité garantit l'existence d'un minimum sur tout compact non vide, tandis que la coercivité permet de ramener la recherche du minimum à un compact suffisamment grand.
\end{methode}
